\section{从稳定系统栈到 Great Transition}

过去十多年，互联网软件形成了一条相对稳定的生产路径：资本购买土地、电力和数据中心，芯片进入机房，云平台把硬件包装成可调用资源，团队在其上训练模型、构建应用，再用安全、合规与治理机制把系统带入现实世界。生成式 AI 没有简单地在这条栈顶增加一个应用类别；它同时抬高了能耗、互连、芯片、数据、模型、产品和治理层的压力，使旧的责任边界与价格结构重新谈判。

\begin{figure}[H]
\centering
\includegraphics[width=0.94\textwidth]{images/00-10-00-course-in-a-nutshell.jpg}
\caption{课程用一条从资本、能源、数据中心到模型、应用与治理的全栈链路组织十周内容。}
\figsource{Stanford Online 官方视频 00:10:00，`images/00-10-00-course-in-a-nutshell.jpg`。}
\end{figure}

\begin{knowledgebox}{读图：先沿依赖方向，再沿价值方向}
从下往上看，资本、土地、电力、芯片和数据中心为计算提供物理条件，cloud/software infrastructure 把异构硬件变成可编排资源，模型与 agent 消耗这些资源形成能力，应用把能力嵌入工作流，治理层处理安全、责任与社会接受度。从上往下看，应用收入和公共价值又决定资本愿意为哪类算力、模型和设施继续投入。两条方向同时存在，所以任何单层“突破”都可能被另一层卡住。
\end{knowledgebox}

\begin{table}[H]
\centering
\begin{tabular}{L{0.18\textwidth}L{0.34\textwidth}L{0.38\textwidth}}
\toprule
层级 & 主要资源或接口 & Frontier Systems 的审计问题 \\
\midrule
资本与能源 & 融资、土地、电网、长期购电 & 投资周期能否匹配技术与需求的不确定性 \\
数据中心与芯片 & 机房、加速器、网络、散热 & 哪些硬件真正可互换，哪里存在锁定 \\
云与软件基础设施 & 调度、虚拟化、存储、可观测性 & 异构资源怎样变成稳定服务接口 \\
模型与 agent & 权重、数据、环境、反馈 & 能力来自静态参数还是持续闭环 \\
应用与部署 & 用户、工作流、分发、收入 & 产品怎样获得高质量 context 与验证信号 \\
安全与治理 & 规范、审计、责任、制度 & 谁定义边界，谁承担失败成本 \\
\bottomrule
\end{tabular}
\caption{课程栈不是公司名单，而是一张资源依赖与责任分配表。}
\end{table}

\subsection{讲者披露与历史位置}

为了理解一门产业课程的证据强弱，本节先把讲者的参与时间线放回图中。早期投资、董事会参与、产品合作或公司创建会提供难以从公开材料获得的细节，也可能让某些问题得到更多注意。好的学习方法不是因为存在利益关系就丢弃全部观点，而是区分第一手观察、公司叙事、宏观推断和价值判断，分别寻找可核验的证据。

\begin{figure}[H]
\centering
\includegraphics[width=0.94\textwidth]{images/00-16-45-background-disclosures.jpg}
\caption{讲者对自身投资、创建或参与项目的时间线披露。}
\figsource{Stanford Online 官方视频 00:16:45，`images/00-16-45-background-disclosures.jpg`。}
\end{figure}

\begin{warningbox}{读图：披露不是装饰性免责声明}
时间线告诉读者哪些判断可能来自内部经验，哪些只能是外部类比。比如，一位长期参与模型公司和算力市场的人更容易看见 capital allocation 与 deployment feedback，却未必代表所有研究团队、公共机构或中小开发者。后文引用讲者关于 context wars、GPU 价格和制度安排的观点时，应把它们标成有来源的经验判断，而不是无主体的客观定律。
\end{warningbox}

\subsection{四类能力约束}

有了全栈地图以后，课程把 AI 进展的限制归为四组：context，compute，capital，culture。这个排序值得重视。社会讨论常把 compute 当作唯一稀缺品，但计算只有在接触正确任务、环境和反馈时才会转化为能力；资本决定扩张速度与风险承担；文化则影响团队愿意测试什么、承认什么失败、把哪些价值写入产品和制度。

\begin{figure}[H]
\centering
\includegraphics[width=0.9\textwidth]{images/00-17-45-capability-bottlenecks.jpg}
\caption{AI 能力进展的四类约束：context、compute、capital 与 culture。}
\figsource{Stanford Online 官方视频 00:17:45，`images/00-17-45-capability-bottlenecks.jpg`。}
\end{figure}

\begin{importantbox}{四类约束如何相互作用}
\term{Context} 决定系统看见哪些任务、工具与反馈；\term{compute} 提供训练和推理预算；\term{capital} 把未来收益预期转成机房、芯片与人才投入；\term{culture} 决定组织如何分享错误、处理安全、选择目标并协调合作。某一层暂时充足时，压力会显现在其他层。诊断系统时要找当前边际约束，不能永远复用同一个答案。
\end{importantbox}

\begin{knowledgebox}{术语消化：bottleneck 不是永久标签}
这里的 bottleneck 指当前增加一单位资源时，最限制整体产出的环节。它会随规模、任务、地区和时间变化。例如，小团队可能先缺 capital，大实验室可能先缺电力与互连，成熟 coding agent 可能更缺可执行环境和真实用户反馈。课程后面频繁跨层，正是因为约束会移动，也会相互放大。
\end{knowledgebox}

\subsection{本章小结}

Great Transition 的含义是整条技术—经济栈正在一起变化。全栈地图提供依赖关系，披露时间线帮助判断观点来源，四类约束则提供动态诊断框架。接下来聚焦第一类约束：context 怎样通过环境、反馈和强化学习形成可持续的能力闭环。

